\chapter{Extinction and Component Topology}

This chapter proves finite extinction for the controlled surgery flow.  The
argument first eliminates components with nontrivial second homotopy, then
uses a loop-space width attached to a nonzero third-homotopy class.

\begin{definition}[Path and ancestry of components]
\label{def:path-of-components}
\dcref{MT:ch18:18.1}
\source{morgan-tian}
\uses{def:controlled-ricci-flow-with-surgery}
\notready
Let $I$ be an interval in the time domain of a controlled Ricci flow with
surgery $({\mathcal M},G)$, with either endpoint allowed to be open, closed,
finite, or infinite.  A \emph{path of components} on $I$ is a connected open
subset ${\mathcal X}\subset\mathbf t^{-1}(I)$ such that
${\mathcal X}(t)={\mathcal X}\cap M_t$ is a connected component of $M_t$ for
every $t\in I$.  On a surgery-free subinterval these components are identified
by the time vector field.  At a singular time $T\in I$, the path ends if its
presurgery component is discarded; otherwise ${\mathcal X}(T)$ is one chosen
continuing postsurgery component.  For $s\le t$, a component
$X_s\subset M_s$ is an \emph{ancestor} of $X_t\subset M_t$, and $X_t$ a
\emph{descendant} of $X_s$, when some path of components on $[s,t]$ has those
endpoint slices.
\end{definition}

\section{Two-Sphere Width Control}

For a $C^1$ map $f\colon S^2\to X$ into a Riemannian manifold, its area is
the integral of the two-dimensional Jacobian, counted with multiplicity.

\begin{definition}[Least nontrivial two-sphere area]
\label{def:two-sphere-width}
\dcref{MT:ch18:18.10}
\source{morgan-tian}
\uses{def:path-of-components}
\notready
Let $(X,g)$ be compact with $\pi_2(X)\ne0$.  Its \emph{two-sphere width} is
\[
 W_2(X,g)=\inf\{\operatorname{Area}_g(f):f\colon S^2\to X\text{ is $C^1$ and not
 null-homotopic}\}.
\]
For a component path $\mathcal X$ in a Ricci flow with surgery, write
\[
 W_2(t)=W_2(\mathcal X(t),g(t))
\]
whenever
$\pi_2(\mathcal X(t))\ne0$.
\end{definition}

\begin{theorem}[Imported Sacks--Uhlenbeck realization of the two-sphere width]
\label{thm:two-sphere-width-realization}
\dcref{MT:ch18:18.10}
\source{morgan-tian}
\uses{def:two-sphere-width}
\notready
If $(X,g)$ is compact and $\pi_2(X)\ne0$, then $W_2(X,g)>0$ and there is a
non-null-homotopic conformal harmonic map $f\colon S^2\to X$, a branched
minimal immersion away from finitely many points, such that
\[
 \operatorname{Area}_g(f)=W_2(X,g).
\]
\end{theorem}

\begin{proof}
This is the Sacks--Uhlenbeck realization used in Morgan--Tian,
Chapter~18, Theorem~18.10.  One minimizes the perturbed energies
$E_\alpha$ over the entire set of maps $S^2\to X$ which are not
null-homotopic, rather than fixing a homotopy class.  The Sacks--Uhlenbeck
energy gap prevents loss of the nontrivial class in a minimizing sequence.
Bubble-tree compactness and the removable-singularity theorem produce a
finite collection of conformal harmonic bubbles; the cut-and-paste argument
in the cited proof shows that at least one bubble is non-null-homotopic and
has area no larger than the original infimum.  Passing \(\alpha\downarrow1\)
therefore gives a non-null-homotopic branched minimal sphere whose area is
$W_2(X,g)$.  The same energy gap gives strict positivity of that infimum.
These compactness and removability statements are the explicit external
Sacks--Uhlenbeck input; no fixed-class or geodesic-cover assumption is being
added.
\end{proof}


\begin{lemma}[Area variation of a minimal two-sphere]
\label{lem:minimal-two-sphere-area-variation}
\dcref{MT:ch18:18.12}
\source{morgan-tian}
\notready
Let $g(t)$ be a Ricci flow on a closed three-manifold near $t_0$, and let
$f\colon S^2\to(M,g(t_0))$ be a branched minimal immersion.  If
$R_{\min}(t_0)=\min_M R_{g(t_0)}$, then
\[
 \left.\frac d{dt}\right|_{t=t_0}\operatorname{Area}_{g(t)}(f)
 \le -4\pi-\frac12R_{\min}(t_0)\operatorname{Area}_{g(t_0)}(f).
\]
\end{lemma}

\begin{proof}
Away from the finite branch set, differentiation of the induced area form
under $\partial_tg=-2\operatorname{Ric}$ gives
\[
 \frac d{dt}\operatorname{Area}_{g(t)}(f)
 =-\int_{S^2}\bigl(R-\operatorname{Ric}(\nu,\nu)\bigr)\,da.
\]
For a minimal surface, the Gauss equation rewrites the integrand as
\[
 -\int_{S^2}K_f\,da-
 \frac12\int_{S^2}(|A|^2+R)\,da.
\]
If the branch orders are $b_1,\ldots,b_q\ge1$, branched Gauss--Bonnet gives
$\int K_f\,da=4\pi+2\pi\sum_jb_j$.  The branch and second-fundamental-form
terms therefore have nonpositive sign.  Finally
$\int R\,da\ge R_{\min}(t_0)\operatorname{Area}_{g(t_0)}(f)$, which proves the claim.
\end{proof}

\begin{lemma}[Fixed-map area distortion]
\label{lem:fixed-sphere-area-distortion}
\dcref{MT:ch18:18.13}
\source{morgan-tian}
\notready
Let $g(t)$ be a Ricci flow on a compact manifold for
$t\in[t_0,t_1]$, suppose $|\operatorname{Ric}_{g(t)}|\le D$, and let
$f\colon S^2\to M$ be $C^1$.  Then, for $s,t\in[t_0,t_1]$,
\[
 e^{-4D|t-s|}\operatorname{Area}_{g(s)}(f)
 \le \operatorname{Area}_{g(t)}(f)
 \le e^{4D|t-s|}\operatorname{Area}_{g(s)}(f).
\]
\end{lemma}

\begin{proof}
Put $A(t)=\operatorname{Area}_{g(t)}(f)$.  At almost every regular point of $f$, the
variation of the two-dimensional Jacobian is one half of the trace of
$\partial_tg=-2\operatorname{Ric}$ on the image plane.  The source norm convention gives
$|A'(t)|\le4DA(t)$.  Thus
$e^{-4D(t-s)}A(t)$ is nonincreasing for $t\ge s$, which proves the upper
bound.  Interchanging $s$ and $t$ proves the lower bound.  Approximation by
smooth maps and dominated convergence remove the regularity restriction.
\end{proof}

\begin{lemma}[Continuity of the two-sphere width at regular times]
\label{lem:two-sphere-width-regular-continuity}
\dcref{MT:ch18:18.11,MT:ch18:18.13}
\source{morgan-tian}
\uses{def:two-sphere-width,thm:two-sphere-width-realization,
  lem:fixed-sphere-area-distortion}
\notready
On a closed surgery-free interval on which the component is identified by
the Ricci-flow time vector field and has nonzero $\pi_2$, the function $W_2$
is continuous.
\end{lemma}

\begin{proof}
Fix $t$ and let $t_i\to t$.  A minimizing sphere at $t$, supplied by
\ref{thm:two-sphere-width-realization}, represents the same nonzero class at
nearby times.  The upper estimate in
\ref{lem:fixed-sphere-area-distortion} gives
$\limsup_iW_2(t_i)\le W_2(t)$.  Conversely, choose a minimizing sphere at
each $t_i$ and view it at time $t$.  The reverse estimate, whose constant is
uniform for $t_i$ near $t$, gives
$W_2(t)\le\liminf_iW_2(t_i)$.  The two inequalities prove continuity.
\end{proof}

\begin{lemma}[Topological comparison for null-sphere surgery]
\label{lem:null-sphere-surgery-topological-comparison}
\dcref{MT:ch18:18.19,MT:ch18:18.22,MT:ch15:15.12}
\source{morgan-tian}
\uses{def:controlled-ricci-flow-with-surgery,
  thm:separating-surgery-map,
  lem:nullhomotopic-embedded-sphere-separates}
\notready
Let $T$ be a surgery time and let $X^-$ be a connected presurgery
component just before $T$.  Assume that every cutting sphere in $X^-$ is
null-homotopic (and hence separating).  Let $X^+$ range over the connected
postsurgery components obtained from $X^-$.  The radial extension of the
flow identification defines, for every sufficiently close regular
$t<T$, a map from $X(t)$ to each selected $X^+$.  Topologically, exactly one
of the following alternatives holds:
\begin{enumerate}
\item if $\pi _2(X^-)$ is nonzero, there is a unique $X^+$ with nonzero
      $\pi _2$, and the corresponding pinch map is a homotopy equivalence;
\item if $\pi _2(X^-)=0$, every $X^+$ has trivial $\pi _2$ and the
      continuing component is either homotopy equivalent to $X^-$ or is a
      homotopy $3$-sphere.
\end{enumerate}
In the second case, if $\pi _3(X^-)\cong\mathbb Z$, the comparison on the
continuing component is either a homotopy equivalence or a degree-one map
and therefore induces an injection on $\pi _3$; the target group is again
$\mathbb Z$.
\end{lemma}

\begin{proof}
The separating-sphere map and the source's one-surgery topology contract
identify the presurgery component with the result of capping the finitely
many cut spheres and recording the discarded sides as connected-sum summands.
When all spheres are null-homotopic, the Morgan--Tian comparison in
Chapter~18 extends the inverse flow map across each cap.  Its cap extensions
are homotopy equivalences on the unique summand carrying a nonzero
$\pi _2$ class; every other summand is a homotopy $3$-sphere.  If the old
$\pi _2$ vanishes, the continuing summand is homotopy equivalent to the old
component or is itself a homotopy $3$-sphere.  In either case the comparison
has degree one whenever it is not a homotopy equivalence, so it injects
$\pi _3$.  This is the topological pinch/degree argument in the cited
source; the metric Lipschitz estimate is supplied separately by
\ref{thm:separating-surgery-map}.
\end{proof}

\begin{theorem}[Component map for a trivial surgery]
\label{thm:trivial-surgery-component-map}
\dcref{MT:ch18:18.18,MT:ch18:18.22,MT:ch15:15.12}
\source{morgan-tian}
\uses{thm:separating-surgery-map,lem:global-surgery-map-lipschitz,
  lem:null-sphere-surgery-topological-comparison,
  lem:nullhomotopic-embedded-sphere-separates}
\notready
Let $T$ be a surgery time, let $X(t)$ be a presurgery component for $t<T$
close to $T$, and suppose every cutting sphere in this component is
null-homotopic (hence separating by
\ref{lem:nullhomotopic-embedded-sphere-separates}).  For each selected
postsurgery component $Y$, for every
$\eta>0$ there is, for all
$t<T$ sufficiently close to $T$, a map
\[
 F_t\colon X(t)\longrightarrow Y,
 \qquad \operatorname{Lip}(F_t)\le1+\eta,
\]
homotopic to the corresponding connected-sum pinch map.  The component with
nonzero $\pi_2$, when it exists, is unique and its pinch map is a homotopy
equivalence.  If $\pi_2(X(t))=0$, then $\pi_2(Y)=0$.  If in addition
$\pi_3(X(t))\cong\mathbb Z$, then $\pi_3(Y)\cong\mathbb Z$ and $(F_t)_*$ is
multiplication by a nonzero integer; in particular it is injective.
\end{theorem}

\begin{proof}
By \ref{lem:nullhomotopic-embedded-sphere-separates}, every cutting sphere is
separating, so the removed-region maps of
\ref{thm:separating-surgery-map} apply simultaneously to the finitely many
necks.  The connected-sum and capping assertions, including the choice of a
continuing component, are supplied by
  \ref{lem:null-sphere-surgery-topological-comparison}.  The topological
  alternatives and the injectivity assertion are exactly
  those of \ref{lem:null-sphere-surgery-topological-comparison}.  Combining
  that comparison with the compact-core extension in
  \ref{lem:global-surgery-map-lipschitz} gives the required map $F_t$ and its
  Lipschitz bound.  Thus the homotopy conclusions do not rely on the metric
  estimate alone, and the global bound is used only with its compact-core and
  finite-cap hypotheses.
\end{proof}

\begin{lemma}[Two-sphere width across a trivial surgery]
\label{lem:two-sphere-width-surgery-control}
\dcref{MT:ch18:18.11,MT:ch15:15.12}
\source{morgan-tian}
\uses{def:path-of-components,def:two-sphere-width,
  thm:trivial-surgery-component-map,lem:nullhomotopic-embedded-sphere-separates}
\notready
Let $T$ be a surgery time after which the surgery spheres in a continuing
component are null-homotopic (and hence separating by
\ref{lem:nullhomotopic-embedded-sphere-separates}), and suppose the selected
postsurgery component
has nonzero $\pi_2$.  Then
\[
 W_2(T)\le\liminf_{t\uparrow T}W_2(t).
\]
\end{lemma}

\begin{proof}
By \ref{thm:trivial-surgery-component-map}, the postsurgery component with
nonzero $\pi_2$ is unique and its pinch map is a homotopy equivalence.  For
every $\eta>0$ and all
$t<T$ sufficiently close to $T$, the pinch equivalence is represented by a
map $F_t\colon\mathcal X(t)\to\mathcal X(T)$ with Lipschitz constant at most
$1+\eta$.  Composition with $F_t$ preserves nontriviality of
a sphere class and multiplies area by at most $(1+\eta)^2$.  Hence
$W_2(T)\le(1+\eta)^2W_2(t)$.  Take a sequence realizing the left liminf and
then let $\eta\downarrow0$.
\end{proof}

\begin{lemma}[Forward inequality for the two-sphere width]
\label{lem:two-sphere-width-forward-inequality}
\dcref{MT:ch18:18.11}
\source{morgan-tian}
\uses{def:path-of-components,def:two-sphere-width,
  thm:two-sphere-width-realization,
  lem:minimal-two-sphere-area-variation,
  lem:two-sphere-width-regular-continuity,
  lem:two-sphere-width-surgery-control}
\notready
Along a component path with nonzero $\pi_2$, after all essential sphere
surgeries have ceased, $W_2$ is continuous at regular times, satisfies
\[
 D^+W_2(t)\le-4\pi-\frac12R_{\min}(t)W_2(t)
\]
there, and obeys the left lower-semicontinuity inequality of
\ref{lem:two-sphere-width-surgery-control} at every surgery time.
\end{lemma}

\begin{proof}
The two continuity assertions are
\ref{lem:two-sphere-width-regular-continuity} and
\ref{lem:two-sphere-width-surgery-control}.  At a regular time $t$, choose a
minimizer $f_t$ from \ref{thm:two-sphere-width-realization}.  Keeping this
map fixed for $h>0$ gives
\[
 W_2(t+h)-W_2(t)
 \le\operatorname{Area}_{g(t+h)}(f_t)-\operatorname{Area}_{g(t)}(f_t).
\]
Divide by $h$, take the upper limit as $h\downarrow0$, and apply
\ref{lem:minimal-two-sphere-area-variation}.
\end{proof}

\begin{lemma}[Finite bound for essential sphere families]
\label{lem:essential-sphere-family-finiteness}
\dcref{MT:ch18:18.8}
\source{morgan-tian}
\notready
For every compact (possibly disconnected) $3$-manifold $X$ there is an
integer $N(X)$ such that every disjoint family of embedded $2$-spheres in
$X$ whose members are individually non-null-homotopic and pairwise
non-homotopic has at most $N(X)$ members.
\end{lemma}

\begin{proof}
It is enough to treat one connected component.  Cutting along a family
$F$ and applying van Kampen gives a graph-of-groups decomposition of
$\pi_1(X)$ with trivial edge groups.  Essentiality forces every vertex of
valence one or two to carry a nontrivial vertex group.  If $k$ is the first
Betti number of the underlying graph, then $k\le\operatorname{rank}H_1(X)$,
and the elementary graph count in the source gives
\[
 2V_1+V_2\ge |F|+3(1-k).
\]
The associated free-product decomposition therefore has at least
$|F|+3(1-k)$ nontrivial factors.  The finitely presented group
$\pi_1(X)$ has only finitely many such factors by the
Kneser--Grushko group bound proved in Morgan--Tian, Chapter~18
(energy1.tex, lines~377--407).  This bounds $|F|$, and taking the
maximum over the finitely many components proves the claim.
\end{proof}

\begin{definition}[Essential sphere family and surgery complexity]
\label{def:essential-sphere-family-complexity}
\dcref{MT:ch18:18.8}
\source{morgan-tian}
\uses{def:controlled-ricci-flow-with-surgery,
  lem:essential-sphere-family-finiteness}
\notready
For a compact three-manifold $X$, an essential sphere family is a finite
disjoint family $F$ of embedded two-spheres such that each member is not
null-homotopic in $X$ and no two distinct members are homotopic in $X$.
Write
\[
 s(X)=\max\{\#F:F\text{ is an essential sphere family in }X\}.
\]
The finiteness of this maximum is supplied by
\ref{lem:essential-sphere-family-finiteness}.
\end{definition}

\begin{lemma}[Essential families survive an essential sphere surgery]
\label{lem:essential-family-survives-essential-surgery}
\dcref{MT:ch18:18.8}
\source{morgan-tian}
\notready
Let $M^-$ be a compact $3$-manifold and let $S_0\subset M^-$ be an embedded
homotopically essential $2$-sphere.  Let $M^+$ be obtained by the standard
surgery on $S_0$, with its two $3$-ball caps, and let
$F^+$ be an essential sphere family in $M^+$ disjoint from the cap balls.
After pushing $F^+$ off the caps and identifying the complement with the
presurgery complement, the family
$F^-:=F^+\cup\{S_0\}$ is essential in $M^-$.
Consequently $s(M^-)>s(M^+)$.
\end{lemma}

\begin{proof}
Push each member of $F^+$ through the cap collars to obtain disjoint
presurgery spheres.  If one of them bounded a homotopy $3$-ball in $M^-$,
the ball either missed $S_0$, in which case the same null-homotopy exists
in $M^+$, or contained $S_0$, in which case $S_0$ would itself be
null-homotopic.  Both alternatives are impossible.  If two members were
homotopic, the intervening $S^2\times I$ either missed $S_0$, giving the
same homotopy in $M^+$, or contained it.  In the latter case a sphere in
$S^2\times I$ is either null-homotopic or parallel to a boundary
component; the first contradicts essentiality of $S_0$, while the second
makes both boundary spheres null-homotopic.  A homotopy involving $S_0$
would likewise make the other sphere null-homotopic after capping.  Thus
$F^-$ is essential, exactly by the three cases in Morgan--Tian,
Chapter~18 (energy1.tex, lines~462--506).
\end{proof}

\begin{theorem}[Finiteness of essential sphere surgeries]
\label{thm:finite-essential-sphere-surgeries}
\dcref{MT:ch18:18.8}
\source{morgan-tian}
\uses{def:controlled-ricci-flow-with-surgery,
  def:essential-sphere-family-complexity,
  lem:essential-sphere-family-finiteness,
  lem:essential-family-survives-essential-surgery,
  thm:finite-interval-surgery-bound}
\notready
In a controlled Ricci flow with surgery, only finitely many
surgery times use a sphere which is non-null-homotopic in its presurgery
component.
\end{theorem}

\begin{proof}
For a surgery time $T$, the source's elementary one-surgery operation
identifies the postsurgery slice with the result of finitely many sphere
surgeries and component removals.  A
null-homotopic surgery leaves $s$ unchanged.  If one of the cutting spheres
$S_0$ is essential, take an essential family $F_T$ realizing $s(M_T)$ and
apply \ref{lem:essential-family-survives-essential-surgery} to the
corresponding elementary surgery.  Thus $s(M_{T^-})>s(M_T)$.  Component
removal cannot increase
$s$.  Thus $s$ is nonincreasing and drops strictly at every time containing
an essential cutting sphere.  The finite upper bound in
\ref{def:essential-sphere-family-complexity}, together with the finite number
of components on each compact slice and the finite-interval surgery bound,
gives finiteness of essential surgery times.
\end{proof}

\section{Eventual Vanishing of \texorpdfstring{$\pi_2$}{pi2}}

\begin{lemma}[Scalar-curvature lower bound]
\label{lem:extinction-scalar-curvature-lower-bound}
\dcref{MT:ch18:18.1,MT:ch15:15.9}
\source{morgan-tian}
\uses{thm:all-time-controlled-surgery-flow}
\notready
For the normalized all-time controlled surgery flow, including the
postsurgery time slices,
\[
 R_{\min}(t)\ge-\frac6{1+4t}.
\]
\end{lemma}

\begin{proof}
The pinching-toward-positive conclusion in
\ref{thm:all-time-controlled-surgery-flow} includes the scalar inequality at
regular times.  At a surgery time the inserted cap and retained neck satisfy
the same pinching inequality, while discarded components do not affect the
minimum on a surviving component.  Thus the estimate passes through every
surgery and holds for all time slices.
\end{proof}

\begin{theorem}[Noncompact two-connected covers are contractible]
\label{thm:noncompact-two-connected-cover-contractible}
\dcref{MT:ch18:18.9}
\source{morgan-tian}
\notready
Let $X$ be a connected three-manifold with $\pi_2(X)=0$.  If its universal
cover $\widetilde X$ is noncompact, then $\widetilde X$ is contractible.
\end{theorem}

\begin{proof}
The cover is simply connected and covering invariance gives
$\pi_2(\widetilde X)=0$.  Degree-two Hurewicz therefore gives
$H_2(\widetilde X;\mathbb Z)=0$, while simple connectivity gives $H_1=0$.
The cover is orientable, and the top homology of a connected noncompact
three-manifold vanishes, so $H_3=0$.  A smooth triangulation has dimension
three, hence $H_i=0$ for $i>3$.

Induct on $n\ge3$.  If $\pi_i(\widetilde X)=0$ for $i<n$, the Hurewicz
theorem identifies $\pi_n(\widetilde X)$ with $H_n(\widetilde X;\mathbb Z)$,
which is zero.  Thus all homotopy groups vanish.  Smooth manifolds and their
covering spaces have CW type, so the map from a point to $\widetilde X$ is a
weak homotopy equivalence between connected CW-type spaces.  Whitehead's
theorem makes it a homotopy equivalence, proving contractibility.
\end{proof}

\begin{theorem}[Eventual vanishing of $\pi_2$]
\label{thm:eventual-pi2-vanishing}
\dcref{MT:ch18:18.9}
\source{morgan-tian}
\uses{def:path-of-components,
  thm:finite-essential-sphere-surgeries,
  thm:trivial-surgery-component-map,
  thm:finite-interval-surgery-bound,
  lem:two-sphere-width-forward-inequality,
  lem:extinction-scalar-curvature-lower-bound,
  lem:forward-dini-comparison-with-downward-jumps,
  thm:noncompact-two-connected-cover-contractible}
\notready
For an all-time controlled surgery flow there is $T_1<\infty$ such that every
component of every time slice $M_T$, $T\ge T_1$, has trivial $\pi_2$.  Any
such component with infinite fundamental group has contractible universal
cover.
\end{theorem}

\begin{proof}
By \ref{thm:finite-essential-sphere-surgeries}, choose $T_0$ after the last
essential sphere surgery.  A component of $M_{T_0}$ with nonzero $\pi_2$ has
a unique continuation with nonzero $\pi_2$ until it is discarded: surgery on
null-homotopic spheres leaves exactly one descendant homotopy equivalent to
the predecessor, while every other descendant is a homotopy three-sphere.

Suppose one such path persists for all $t\ge T_0$.  By
\ref{lem:two-sphere-width-forward-inequality} and
\ref{lem:extinction-scalar-curvature-lower-bound}, its positive width obeys
\[
 D^+W_2(t)\le-4\pi+\frac3{1+4t}W_2(t).
\]
Let $w(T_0)=W_2(T_0)$ and
$w'=-4\pi+3w/(1+4t)$.  The integrating factor
$(1+4t)^{-3/4}$ gives
\[
 w(t)=W_2(T_0)\left(\frac{1+4t}{1+4T_0}\right)^{3/4}
 +4\pi(1+4T_0)^{1/4}(1+4t)^{3/4}-4\pi(1+4t).
\]
Continuity on each regular interval, the surgery-time downward liminf, and
the finite-partition comparison lemma
\ref{lem:forward-dini-comparison-with-downward-jumps} give
$W_2(t)\le w(t)$.  The last term in
$w$ has order $t$, while the first two have order $t^{3/4}$, so $w$ is
eventually negative, contradicting $W_2(t)>0$.

There are finitely many components at $T_0$, hence one finite $T_1$ works for
all their continuations.  Triviality of $\pi_2$ persists under later trivial
sphere surgeries.  The final assertion follows from
\ref{thm:noncompact-two-connected-cover-contractible}, since an infinite
fundamental group makes the universal cover noncompact.
\end{proof}
\section{Late-Component \texorpdfstring{$\pi_3$}{pi3} Topology}

\begin{lemma}[Imported finitely generated Kurosh factor contract]
\label{lem:finitely-generated-kurosh-free-factor}
\dcref{MT:ch18:18.19}
\source{morgan-tian}
\notready
Let
$G=G_1*\cdots*G_m*F_r$, with the $G_i$ finite, and let $H$ be a finitely
generated free factor of $G$.  Then $H$ is a free product of finitely many
finite groups and a finitely generated free group.
\end{lemma}

\begin{proof}
This is the finitely generated form of the Kurosh subgroup theorem, used in
Morgan--Tian, Chapter~18, Theorem~18.19.  Realize the free product as the
fundamental group of its Bass--Serre tree of groups.  A finitely generated
free factor acts on that tree with a finite quotient of its minimal subtree:
only finitely many conjugates of the finite vertex groups occur, and all
remaining vertex stabilizers are trivial.  The quotient graph of groups is
therefore finite, and its fundamental group is a free product of those
finitely many finite stabilizers and a finitely generated free group (the
rank is the first Betti number of the finite quotient graph).  This gives the
stated decomposition, with the finite-generation hypothesis exactly the one
needed to make the quotient graph finite.
\end{proof}

\begin{theorem}[Persistence of the fundamental-group class]
\label{thm:component-fundamental-group-persistence}
\dcref{MT:ch18:18.19,MT:ch15:15.3}
\source{morgan-tian}
\uses{def:path-of-components,def:controlled-ricci-flow-with-surgery,
  thm:trivial-surgery-component-map,
  lem:finitely-generated-kurosh-free-factor,
  thm:finite-interval-surgery-bound}
\notready
Let $[0,T]$ be a bounded interval containing only finitely many surgery times.
Suppose the fundamental group of every initial component of a controlled
surgery flow is a free product of finitely many finite groups and infinite
cyclic groups.  Then the same is true for every component of every time slice
up to $T$.
\end{theorem}

\begin{proof}
The time vector field identifies components diffeomorphically on every
regular interval.  At a surgery time, the connected-sum decomposition is the
one-surgery topology contract included in the component-map theorem
\ref{thm:trivial-surgery-component-map}.  Van Kampen therefore
makes the fundamental group of each postsurgery component a free factor of the
corresponding presurgery group; the discarded spherical or $S^2$-bundle
factors are finite or infinite cyclic.  This is compatible with the
connected-sum pinch maps in \ref{thm:trivial-surgery-component-map}.

The group-theoretic closure is exactly
\ref{lem:finitely-generated-kurosh-free-factor}; a free factor is finitely
generated because it is a retract of the finitely generated presurgery group.
The finite-interval surgery bound makes the induction over surgery times
finite on every bounded interval.
\end{proof}

\begin{lemma}[Hempel essential-sphere input]
\label{lem:hempel-essential-sphere-input}
\dcref{MT:ch18:18.20}
\source{morgan-tian}
\notready
Let $X$ be a compact connected boundaryless $3$-manifold.  If
$\pi_1(X)$ is a nontrivial free product, or if $\pi_1(X)\cong\mathbb Z$, then
$X$ contains an embedded $2$-sphere which is not null-homotopic.  In
particular, $\pi_2(X)\ne0$.
\end{lemma}

\begin{proof}
This is the exact imported topology statement used by Morgan--Tian
Chapter~18, Theorem~18.20: their cited Hempel Theorem~5.2 handles the
infinite-cyclic case and Hempel Theorem~7.1 handles a nontrivial free-product
decomposition.  The source's ``compact $3$-manifold'' is interpreted in the
boundaryless time-slice convention displayed in the statement.  The imported
conclusion is an embedded essential sphere, not merely a formal
sphere-system component, so it directly supplies the nonzero $\pi_2$ class.
\end{proof}

\begin{theorem}[Imported Morgan--Tian/Hempel free-product obstruction to vanishing $\pi_2$]
\label{thm:free-product-pi2-obstruction}
\dcref{MT:ch18:18.20}
\source{morgan-tian}
\uses{lem:hempel-essential-sphere-input}
\notready
If $X$ is a closed connected three-manifold with $\pi_1(X)$ a nontrivial
free product or $\pi_1(X)\cong\mathbb Z$, then $\pi_2(X)\ne0$.
\end{theorem}

\begin{proof}
Apply \ref{lem:hempel-essential-sphere-input}.  Its essential sphere is not
null-homotopic, so its based homotopy class is nonzero in $\pi_2(X)$; no
maximal-sphere-system or target-equivalent decomposition is being used as an
unstated premise.
\end{proof}


\begin{theorem}[Finite fundamental group and the universal cover]
\label{thm:finite-cover-homotopy-three-sphere}
\dcref{MT:ch18:18.19}
\source{morgan-tian}
\notready
Let $X$ be a closed connected three-manifold with finite fundamental group
and $\pi_2(X)=0$.  Its universal cover $\widetilde X$ is a homotopy
three-sphere, and covering projection induces
\[
 \pi_3(\widetilde X)\cong\pi_3(X)\cong\mathbb Z.
\]
\end{theorem}

\begin{proof}
The universal cover is finite-sheeted, hence compact, closed, and simply
connected.  It is orientable, and covering invariance gives
$\pi_2(\widetilde X)=0$.  Poincare duality gives
$H_3(\widetilde X;\mathbb Z)\cong\mathbb Z$, while degree-two Hurewicz gives
$H_2=0$ and simple connectivity gives $H_1=0$.  Degree-three Hurewicz now
identifies $\pi_3(\widetilde X)$ with $H_3$.

Choose a map $S^3\to\widetilde X$ representing a generator.  It is an
isomorphism on every integral homology group.  Both spaces are simply
connected CW-type spaces, so the homological Whitehead theorem makes this map
a homotopy equivalence.  Thus $\widetilde X$ is a homotopy three-sphere.
Covering invariance in degrees at least two gives the asserted isomorphism
with $\pi_3(X)$.
\end{proof}

\begin{theorem}[Nonzero $\pi_3$ on a late surviving component]
\label{thm:late-component-nontrivial-pi3}
\dcref{MT:ch18:18.19}
\source{morgan-tian}
\uses{def:path-of-components,thm:eventual-pi2-vanishing,
  thm:component-fundamental-group-persistence,
  thm:free-product-pi2-obstruction,
  thm:finite-cover-homotopy-three-sphere}
\notready
Let $T_1$ be as in \ref{thm:eventual-pi2-vanishing}, and assume the initial
fundamental groups are free products of finite groups and infinite cyclic
groups.  If $\mathcal X$ is a component path defined on $[T_1,T_2]$, then
\[
 \pi_3(\mathcal X(T_1))\cong\mathbb Z.
\]
In particular it contains a nonzero class.
\end{theorem}

\begin{proof}
Put $X_1=\mathcal X(T_1)$.  By \ref{thm:eventual-pi2-vanishing},
$\pi_2(X_1)=0$.  The persistence theorem
\ref{thm:component-fundamental-group-persistence} expresses $\pi_1(X_1)$
as a free product of finite and infinite cyclic
groups.  If more than one nontrivial factor occurred, or if an infinite
cyclic factor occurred, \ref{thm:free-product-pi2-obstruction} would
contradict the vanishing of $\pi_2$.  Thus the fundamental group is finite,
possibly trivial.  Apply \ref{thm:finite-cover-homotopy-three-sphere}.
\end{proof}

\section{The Loop-Space Width Interface}

Let $\Lambda_0X$ denote the component of null-homotopic loops in the free
$C^1$ loop space of a compact manifold $X$, and let $*$ be the constant loop
at a chosen base point $x_0$.

\begin{lemma}[The quotient and wedge model]
\label{lem:loop-space-quotient-wedge}
\dcref{MT:ch18:18.16}
\source{morgan-tian}
\notready
For $c_0\in S^2$, the quotient
\[
 Q=(S^2\times S^1)/(\{c_0\}\times S^1)
\]
is homotopy equivalent to $S^2\vee S^3$.  Consequently, for every based
space $X$ of CW type,
\[
 [Q,X]_*\cong\pi_2(X)\times\pi_3(X).
\]
\end{lemma}

\begin{proof}
Give $S^2$ one zero-cell and one two-cell, choosing the zero-cell to be
$c_0$, and give $S^1$ one zero-cell and one one-cell.  The product has cells
in dimensions $0,1,2,3$.  Collapsing $\{c_0\}\times S^1$ removes the
one-cell.  The attaching map of the remaining three-cell is the Whitehead
product of the two-cell with the collapsed one-cell, hence is null after the
collapse.  The quotient CW complex is therefore a two-sphere with a
three-cell attached at the base point, namely $S^2\vee S^3$.  Based maps out
of a wedge are pairs of based maps, which gives the final identification.
\end{proof}

\begin{lemma}[Lifting null-loop families]
\label{lem:loop-space-universal-cover-lift}
\dcref{MT:ch18:18.16}
\source{morgan-tian}
\notready
Let $p\colon\widetilde X\to X$ be the universal cover of a connected
manifold.  A continuous family $\Gamma\colon S^2\to\Lambda_0X$ lifts, after
choosing one lift of one marked point, to a continuous family of loops in
$\widetilde X$.  Homotopies lift in the same way, and changing the initial
lift acts by one deck transformation on the entire lifted family.
\end{lemma}

\begin{proof}
The adjoint map $F\colon S^2\times S^1\to X$ sends every circle
$\{c\}\times S^1$ to a null-homotopic loop.  The induced homomorphism on
$\pi_1(S^2\times S^1)\cong\mathbb Z$ is therefore zero, so the covering
lifting criterion gives a lift $\widetilde F$.  Uniqueness after fixing one
point proves continuity of the family and the assertion about deck
transformations.  Apply the same argument to
$S^2\times S^1\times[0,1]$ for a homotopy.
\end{proof}

\begin{theorem}[Free-loop-space identification of $\pi_3$]
\label{thm:loop-space-pi3-identification}
\dcref{MT:ch18:18.16}
\source{morgan-tian}
\uses{lem:loop-space-quotient-wedge,
  lem:loop-space-universal-cover-lift}
\notready
If $X$ is compact, connected, and $\pi_2(X,x_0)=0$, then adjunction induces
an isomorphism
\[
 \pi_2(\Lambda_0X,*)\cong\pi_3(X,x_0).
\]
Free homotopy classes of not necessarily based maps
$S^2\to\Lambda_0X$ correspond to the orbit set of $\pi_3(X,x_0)$ under its
natural $\pi_1(X,x_0)$-action.
\end{theorem}

\begin{proof}
For based families, compact-open adjunction identifies a map
$S^2\to\Lambda_0X$ with a map $S^2\times S^1\to X$ which collapses
$\{c_0\}\times S^1$.  By \ref{lem:loop-space-quotient-wedge}, its based
homotopy class is a pair in $\pi_2(X)\times\pi_3(X)$.  The first factor
vanishes, giving the displayed isomorphism.  Passing between continuous and
$C^1$ loops does not change the homotopy type: smoothing in finitely many
charts, relative to already smooth loops, gives mutually inverse homotopies.

For an unbased family in $\Lambda_0X$, use
\ref{lem:loop-space-universal-cover-lift}.  A choice of lift turns it into the
based construction after moving one constant loop to the chosen base point.
A different lift applies a deck transformation throughout and changes the
based representative by the natural $\pi_1(X)$-action.  Conversely, a change
by this action is realized by moving the base point around the corresponding
loop.  Thus forgetting the base point gives precisely the stated orbit set.
\end{proof}

\begin{definition}[Loop-family width]
\label{def:loop-family-width}
\dcref{MT:ch18:18.17}
\source{morgan-tian}
\uses{thm:loop-space-pi3-identification}
\notready
For a null-homotopic $C^1$ loop $\gamma$ in a compact Riemannian manifold,
let
\[
 \mathcal A_g(\gamma)=\inf_F\operatorname{Area}_g(F),
\]
where $F\colon D^2\to X$ ranges over Lipschitz maps whose boundary agrees
with $\gamma$ up to orientation-preserving reparameterization.  For a
continuous family $\Gamma\colon S^2\to\Lambda_0X$, put
\[
 W_g(\Gamma)=\sup_{c\in S^2}\mathcal A_g(\Gamma(c)).
\]
If $\xi\in\pi_3(X)$, let $[\xi]$ be its $\pi_1(X)$-orbit.  By
\ref{thm:loop-space-pi3-identification}, this orbit corresponds to a free
homotopy class of loop families.  Define
\[
 W_g(\xi)=\inf_{[\Gamma]=[\xi]}W_g(\Gamma),
\]
where representatives need not preserve the base point.  Thus $W_g$ is
constant on every $\pi_1(X)$-orbit.
\end{definition}

\begin{lemma}[Continuity of least spanning area]
\label{lem:least-spanning-area-continuity}
\dcref{MT:ch18:18.17}
\source{morgan-tian}
\notready
The function $\gamma\mapsto\mathcal A_g(\gamma)$ is finite and continuous
on $\Lambda_0X$ in the $C^1$ topology.  More precisely, two sufficiently
$C^1$-close loops cobound a Lipschitz annulus whose area tends to zero with
their $C^1$ distance, and
\[
 |\mathcal A_g(\gamma_0)-\mathcal A_g(\gamma_1)|
 \le\inf\{\operatorname{Area}_g(A):\partial A=\gamma_0\sqcup\gamma_1\}.
\]
\end{lemma}

\begin{proof}
A null-homotopy can be approximated by a Lipschitz disk, so the infimum is
finite.  If $\gamma_0$ and $\gamma_1$ are sufficiently $C^1$-close, join
$\gamma_0(s)$ to $\gamma_1(s)$ by the unique short geodesic.  Smooth
dependence of these geodesics produces an annulus whose two derivatives are
uniformly bounded respectively by the pointwise distance and by the common
$C^1$ bound.  Its area therefore tends to zero with the $C^1$ distance.
Gluing this annulus to a disk spanning either boundary gives both directions
of the displayed inequality and hence continuity.
\end{proof}

\section{The Width Comparison ODE}

\begin{definition}[Comparison equation]
\label{def:loop-width-comparison-ode}
\dcref{MT:ch18:18.23}
\source{morgan-tian}
\notready
For a compact Ricci flow $(M,g(t))$, $t\in[t_0,t_1]$, and
$t'\in[t_0,t_1]$, let $w_{a,t'}$ be the solution of
\[
 w'=-2\pi-\frac12R_{\min}(t)w,
 \qquad w(t')=a.
\]
Write $w_a=w_{a,t_0}$.
\end{definition}

\begin{lemma}[Integrating-factor formula]
\label{lem:loop-width-ode-formula}
\dcref{MT:ch18:18.25,MT:ch18:18.26}
\source{morgan-tian}
\uses{def:loop-width-comparison-ode}
\notready
Put
\[
 B_{t'}(t)=\int_{t'}^t\frac12R_{\min}(s)\,ds.
\]
Then, for $t''\ge t'$,
\[
 w_{a,t'}(t'')=e^{-B_{t'}(t'')}
 \left(a-2\pi\int_{t'}^{t''}e^{B_{t'}(s)}\,ds\right).
\]
If $a'\ge a$, then
\[
 w_{a',t'}(t'')-w_{a,t'}(t'')
 =(a'-a)e^{-B_{t'}(t'')}\ge0.
\]
On a fixed compact time interval this difference is bounded by a constant
times $a'-a$.
\end{lemma}

\begin{proof}
Multiplication of the equation in \ref{def:loop-width-comparison-ode} by
$e^{B_{t'}(t)}$ gives
\[
 \frac d{dt}\bigl(e^{B_{t'}(t)}w_{a,t'}(t)\bigr)
 =-2\pi e^{B_{t'}(t)}.
\]
Integration from $t'$ to $t''$ proves the first formula.  Subtracting the
formulas with initial values $a'$ and $a$ proves the second.  Boundedness of
$R_{\min}$ on a compact regular interval bounds the exponential factor.
\end{proof}
\section{Controlled Loop-Family Deformation}

\begin{theorem}[Plateau disk existence and boundary regularity]
\label{thm:plateau-disk-existence-regularity}
\dcref{MT:ch19:19.2}
\source{morgan-tian}
\source{white-classical-area-minimizing-surfaces:page-0001}
\source{white-classical-area-minimizing-surfaces:page-0002}
\notready
Let $(N,g)$ be a compact real-analytic Riemannian three-manifold and let
$\gamma\colon S^1\to N$ be an embedded real-analytic null-homotopic curve.
The infimum of the areas of $W^{1,2}$ disk maps whose boundary trace is
$\gamma$ (up to an orientation-preserving weakly monotone
reparameterization) is attained by a conformal harmonic Douglas disk.  Its
interior branch set is finite, it is real analytic up to the boundary, and it
has no true boundary branch point.
\end{theorem}
\begin{proof}
Morgan--Tian's Plateau input in Chapter~19, Section~19.2, is the Douglas
direct method followed by the Morrey boundary theorem; it supplies attainment
and the conformal-harmonic representative with the displayed trace
convention.  White's registered theorem states that an area-minimizing disk
in a real-analytic Riemannian target has a finite branch set and excludes a
true branch point on a real-analytic boundary (pages~295--296).  White is used
only for that boundary conclusion; false branch points are not excluded.
\end{proof}

\begin{lemma}[Least-disk area is continuous under a boundary isotopy]
\label{lem:least-disk-area-continuity-under-isotopy}
\dcref{MT:ch19:19.4}
\source{morgan-tian}
\uses{thm:plateau-disk-existence-regularity}
\notready
Let $(N,g)$ be compact and let
$\Gamma\colon S^1\times[0,1]\to N$ be a $C^2$ isotopy such that every
$\Gamma_s$ is an embedded null-homotopic curve.  If
$\mathcal A_g(\Gamma_s)$ denotes the least spanning-disk area, then
$s\mapsto\mathcal A_g(\Gamma_s)$ is continuous.
\end{lemma}
\begin{proof}
On a compact parameter interval the isotopy sweeps a $C^2$ collar whose area
over $[s,s+h]$ is $O(|h|)$.  Gluing that collar to a minimizing disk for
$\Gamma_s$ gives
$\mathcal A_g(\Gamma_{s+h})\le\mathcal A_g(\Gamma_s)+o(1)$; reversing the
collar gives the opposite inequality.  The minimizing disks needed at the
endpoints are supplied by the preceding Douglas contract after real-analytic
approximation, and the collar estimate is unchanged in the limit.  This is
the moving-boundary argument of Morgan--Tian Section~19.4.
\end{proof}

\begin{lemma}[Single immersed loop approximation by embedded boundaries]
\label{lem:immersed-loop-embedded-approximation}
\dcref{MT:ch19:19.4}
\source{morgan-tian}
\uses{lem:least-disk-area-continuity-under-isotopy}
\notready
Let $\gamma\colon S^1\to N$ be a $C^2$ immersed null-homotopic loop in a
smooth three-manifold.  For every $\varepsilon>0$ there is an embedded $C^2$
null-homotopic loop $\gamma_\varepsilon$, $C^2$-close to $\gamma$, and a
thin collar homotopy between them such that
\[
 \bigl|\mathcal A_g(\gamma_\varepsilon)-\mathcal A_g(\gamma)\bigr|<\varepsilon.
\]
The assertion is for each loop (or each member of a finite net), not a claim
that an arbitrary positive-dimensional family can be made simultaneously
embedded.
\end{lemma}
\begin{proof}
First make an arbitrarily small generic $C^2$ perturbation, so the only
self-intersections are finitely many transverse double points.  In pairwise
disjoint target balls, separate the two local branches at each such point;
general position in dimension three then gives an embedded curve while
preserving the null homotopy.  The perturbation is joined to the original by a
collar with area tending to zero.  Apply the preceding isotopy-continuity
lemma to the two boundary traces.  This is the generic-boundary
approximation used in Morgan--Tian Section~19.4, stated only in the
single-loop form consumed by the variation argument.
\end{proof}

\begin{lemma}[Plateau disk and moving-boundary contract]
\label{lem:plateau-disk-moving-boundary-contract}
\dcref{MT:ch19:19.2,MT:ch19:19.4}
\source{morgan-tian}
\source{white-classical-area-minimizing-surfaces:page-0001}
\source{white-classical-area-minimizing-surfaces:page-0002}
\uses{def:loop-family-width,thm:plateau-disk-existence-regularity,
  lem:least-disk-area-continuity-under-isotopy,
  lem:immersed-loop-embedded-approximation}
\notready
The three preceding contracts together provide exactly the disk existence,
boundary regularity, isotopy continuity, and single-loop immersed-boundary
approximation used by the least-area variation theorem.  No annular
attainment or simultaneous family-embedding conclusion is included.
\end{lemma}
\begin{proof}
For an embedded analytic boundary use
\ref{thm:plateau-disk-existence-regularity}; use
\ref{lem:least-disk-area-continuity-under-isotopy} for the collar comparison,
and pass from an immersed boundary by
\ref{lem:immersed-loop-embedded-approximation}.  Each appeal has the exact
regularity and topology hypotheses displayed in its statement.
\end{proof}

\begin{theorem}[Least-disk variation interface]
\label{thm:least-disk-area-variation-interface}
\dcref{MT:ch19:19.2,MT:ch19:19.4}
\source{morgan-tian}
\source{white-classical-area-minimizing-surfaces:page-0001}
\source{white-classical-area-minimizing-surfaces:page-0002}
\uses{def:loop-family-width,lem:plateau-disk-moving-boundary-contract}
\notready
Let $(M,g(t))$, $t_0\le t\le t_1$, be a Ricci flow on a compact
three-manifold, and let $\gamma(t)$ be a $C^2$ curve-shrinking flow of
null-homotopic immersed loops.  Then $t\mapsto\mathcal A_{g(t)}(\gamma(t))$
is continuous and, at every time at which the flow is smooth,
\[
 D^+\mathcal A_{g(t)}(\gamma(t))
 \le-2\pi-\frac12R_{\min}(t)
 \mathcal A_{g(t)}(\gamma(t)).
\]
\end{theorem}

\begin{proof}
For an embedded boundary at a regular time, use
\ref{lem:plateau-disk-moving-boundary-contract} to obtain a least disk with
conformal harmonic interior and the required boundary regularity.  Vary the
metric by $\partial_tg=-2\operatorname{Ric}$ and vary the
boundary by the curve-shrinking vector field.  The first variation of the
minimizing disk vanishes in its interior; the boundary term is the geodesic
curvature of the moving boundary.  Gauss--Bonnet on the disk, together with
minimality and the Gauss equation, gives
\[
 -\int_D\operatorname{tr}_{TD}\operatorname{Ric}\,da
 -\int_{\partial D}k_g\,ds
 =-\int_DK_D\,da
  -\frac12\int_D(|A|^2+R)\,da
  -\int_{\partial D}k_g\,ds
 \le -2\pi-\tfrac12R_{\min}(t)\operatorname{Area}_{g(t)}(D).
\]
Taking the infimum over disks yields the displayed upper Dini derivative; the
same first-variation formula and thin annuli swept out by the boundary give
continuity.  An immersed $C^2$ loop is approximated in $C^2$ by embedded loops,
 the corresponding least disks are compared by thin collar annuli, and the
 upper derivative inequality passes to the limit.  The Plateau regularity,
 moving-boundary variation, and this approximation are the imported contract
 recorded in Morgan--Tian 19.2 and 19.4.
\end{proof}

\begin{definition}[Nondegenerate Douglas annulus class]
\label{def:nondegenerate-douglas-annulus-class}
\dcref{MT:ch19:19.15}
\source{morgan-tian}
\uses{def:ramp-curve}
\notready
Let $(M,h)$ be a compact real-analytic Riemannian three-manifold, put
$\bar h=h+ds^2$ on $M\times S^1_\lambda$, and let
$c_0,c_1\colon S^1\to M\times S^1_\lambda$ be disjoint embedded
real-analytic degree-one ramps in one free homotopy class.  For $\rho>0$,
write $\mathscr A_\rho(c_0,c_1)$ for the class of $W^{1,2}$ maps $F$ from a
compact annulus whose boundary traces are $c_0,c_1$ up to
orientation-preserving weakly monotone reparameterizations and such that
every essential simple closed curve $\alpha$ in the domain satisfies
\[
 L_{\bar h}(F\circ\alpha)\ge\rho.
\]
Area is counted with multiplicity.  The class is called nondegenerate when it
is nonempty.
\end{definition}
\begin{proof}
This is a definition.  The explicit trace and length conditions are the
Douglas nondegeneration convention used in Morgan--Tian Proposition~19.15;
the product metric fixes the scale and the meaning of every length below.
\end{proof}

\begin{theorem}[Douglas attainment and regularity for a nondegenerate annulus]
\label{thm:annular-douglas-minimizer}
\dcref{MT:ch19:19.15,
  white-classical-area-minimizing-surfaces:page-0001,
  white-classical-area-minimizing-surfaces:page-0002}
\source{morgan-tian}
\source{white-classical-area-minimizing-surfaces:page-0001}
\source{white-classical-area-minimizing-surfaces:page-0002}
\uses{def:nondegenerate-douglas-annulus-class}
\notready
Let $\rho>0$ and assume that the class
$\mathscr A_\rho(c_0,c_1)$ of
\ref{def:nondegenerate-douglas-annulus-class} is nonempty.  Then its annular
area infimum is attained by a classical conformal harmonic annulus $A$.  Its
interior branch set is finite, it is real analytic up to its real-analytic
boundary, and it has no true boundary branch point.  False boundary branch
points are not excluded.  There are interior parallel boundary curves
avoiding the finite branch set which converge to the original boundary and
whose intervening collar areas tend to zero.
\end{theorem}

\begin{proof}
This is the product/ramp contract in Morgan--Tian, Proposition~19.15
(\texttt{energy1.tex}, lines~1920--1950 and 2985--3018), with the admissible
class made explicit in \ref{def:nondegenerate-douglas-annulus-class}.  The
degree-one projection is the source of the positive essential-loop lower
bound in the applications; the direct Douglas compactness argument in the
cited proposition then attains the annular infimum.  The
proposition imports the Hildebrandt--Morrey regularity conclusion for the
resulting conformal harmonic annulus: finite interior branch set and real
analyticity up to the real-analytic boundary.  White's registered theorem
(pages~295--296) excludes true boundary branch points but not false ones.
Parallel curves avoiding the finite interior branch set give the stated
collar exhaustion with areas tending to zero.
\end{proof}

\paragraph{Annular attainment convention.}
Every later annular consumer invokes the preceding theorem with its explicit
positive essential-loop lower bound.  Thus Douglas attainment is a produced
conclusion of the registered Morgan--Tian source contract, not a hidden
hypothesis.  White is used only for the stated exclusion of true boundary
branch points; false boundary branches remain part of the contract.


\begin{lemma}[Ramp nondegeneration]
\label{lem:ramp-annulus-nondegeneration}
\dcref{MT:ch19:19.15}
\source{morgan-tian}
\uses{def:ramp-curve}
\notready
Let $(M,h)$ be a compact Riemannian three-manifold and equip
$M\times S^1_\lambda$ with the product metric $\bar h=h+ds^2$.  Let
$c_0,c_1\colon S^1\to M\times S^1_\lambda$ be degree-one ramps in the
same free homotopy class, and let
$F\colon S^1\times[0,1]\to M\times S^1_\lambda$ span them.  For every
essential simple closed domain curve $\alpha$,
\[
 L_{\bar h}(F\circ\alpha)\ge L_{ds^2}(S^1_\lambda)>0.
\]
The same bound holds while both curves evolve by smooth curve shrinking.
\end{lemma}

\begin{proof}
An essential simple closed curve in the annulus is freely homotopic, with one
of the two orientations, to either boundary component.  Its composition with
the circle projection therefore has degree $1$ or $-1$.  Such a map onto
$S^1_\lambda$ has length at least the circumference of the target, and product
projection does not increase length.  A smooth curve flow is a homotopy, so
the circle degree and common free homotopy class are preserved.
\end{proof}

\begin{lemma}[Smooth-flow bridge to the annular minimizer contract]
\label{lem:smooth-annular-minimizer-bridge}
\dcref{MT:ch19:19.15}
\source{morgan-tian}
\uses{def:ramp-curve,thm:annular-douglas-minimizer,
  lem:ramp-annulus-nondegeneration}
\notready
Let $(M,g(t))$, $t_0\le t\le t_1$, be a smooth Ricci flow on a compact
three-manifold, and let $c_0(t),c_1(t)$ be smooth degree-one ramp
curve-shrinking flows in $M\times S^1_\lambda$ in one free homotopy class.
At each $t>t_0$ the metric and the two evolved curves are real analytic, the
annular nondegeneration bound of
\ref{lem:ramp-annulus-nondegeneration} holds.  The preceding theorem then
produces a classical Douglas solution with the stated finite branch set and
collar exhaustion.  The same assertion at $t=t_0$ and all right-hand
variation statements are understood as limits of real-analytic approximants
to the initial metric and curves.
After the false boundary branches are reparameterized, the pullback metric is
regularized near its finitely many interior branch points: for every prescribed
area tolerance it is replaced by a smooth metric with at most twice the
annulus area and with a finite ambient-dependent upper Gaussian-curvature
bound.  This regularization is part of the annular comparison contract.
\end{lemma}

\begin{proof}
For positive time, the immediate real-analyticity of a Ricci flow and of a
curve-shrinking solution is the regularization step used in the proof of
Morgan--Tian 19.15.  Apply
\ref{thm:annular-douglas-minimizer}; the positive degree-one lower bound
supplies its nondegeneration input and hence its attainment conclusion.  At
the initial time, approximate the
smooth metric and the two curves in $C^2$ by real-analytic data in the same
free homotopy class.  The degree-one projection bound persists (with a fixed
positive lower bound), so the analytic minimizers have a uniform
nondegeneration constant.  The direct-method compactness and the collar
exhaustion pass to the smooth limit.  The same approximation, followed by
the upper right derivative, passes the first-variation inequality to the
original data.  This is precisely the analytic-regularization paragraph in
the source proof, rather than an extra analyticity hypothesis on the Ricci
flow consumer.
\end{proof}

\begin{definition}[Finite small-annulus net]
\label{def:finite-small-annulus-net}
\dcref{MT:ch18:18.24}
\source{morgan-tian}
\notready
Let $\Gamma\colon S^2\to\Lambda M$ be a continuous $C^1$-loop family.  A
finite set ${\mathcal S}\subset S^2$ is a \emph{$\nu$-annular net} for
$\Gamma$ if, for every $c\in S^2$, there are $\widehat c\in{\mathcal S}$ and
an admissible Lipschitz annulus from $\Gamma(c)$ to $\Gamma(\widehat c)$ of
area less than $\nu$.  Admissible means that the two boundary traces agree
with the displayed loops up to orientation-preserving weakly monotone
reparameterizations, as in the annular-area infimum.  The same terminology is
used for a ramp family in $M\times S^1_\lambda$, with area measured in the
product metric.  In the source the admissible annuli are the swept annuli
associated to oriented parameter arcs; this orientation clause is part of the
net contract.
\end{definition}

\begin{lemma}[Finite annular nets and a common ramp scale]
\label{lem:finite-net-ramp-nondegeneration}
\dcref{MT:ch18:18.24,MT:ch19:19.15}
\source{morgan-tian}
\uses{def:finite-small-annulus-net,
  thm:controlled-curve-shrinking-existence,
  lem:ramp-annulus-nondegeneration}
\notready
Let $g(t)$, $t_0\le t\le t_1$, be a Ricci flow on a compact manifold, and let
$\widetilde\Gamma\colon S^2\to\Lambda_0M$ and
$\widetilde\Gamma_c^\lambda(t)$ be the approximation and ramp flows supplied
by \ref{thm:controlled-curve-shrinking-existence}.  For every $\mu>0$ there
are a finite set ${\mathcal S}\subset S^2$ and $\lambda_{\rm net}>0$ such
that ${\mathcal S}$ is a $\mu/2$-annular net for
$\widetilde\Gamma$, and, for every
$0<\lambda<\lambda_{\rm net}$, it is a $\mu$-annular net for the initial
ramp family $\widetilde\Gamma^\lambda$.  Thus, for each $c$, one can choose
$\widehat c\in{\mathcal S}$ so that the initial annular-area infimum between
$\widetilde\Gamma_c^\lambda(t_0)$ and
$\widetilde\Gamma_{\widehat c}^\lambda(t_0)$ is less than $\mu$.

At every later time the two boundary curves are embedded degree-one ramps in
the same free homotopy class.  Consequently every essential simple closed
curve in every spanning annulus for the pair has length at least
\[
 \rho_\lambda=L(S^1_\lambda)=\lambda,
\]
uniformly over $c$, the finite net, and $[t_0,t_1]$ for this fixed
$\lambda$.  Finally, after ${\mathcal S}$ is chosen, any collection of
positive pointwise thresholds $\lambda_{\widehat c}$,
$\widehat c\in{\mathcal S}$, admits one common scale
\[
 0<\lambda<\min\bigl\{\lambda_{\rm net},
     \lambda_{\widehat c}:\widehat c\in{\mathcal S}\bigr\}.
\]
No lower bound for $\rho_\lambda$ uniform as $\lambda\downarrow0$ is
asserted or needed.
\end{lemma}

\begin{proof}
Continuity of $\widetilde\Gamma$ in the $C^1$ topology gives, around every
$c$, a parameter neighborhood in which any two loops are joined by a short
geodesic ruled annulus of area less than $\mu/2$.  Compactness of $S^2$ gives
a finite subcover and hence a finite $\mu/2$-annular net.  Lift the finitely
many local ruled-annulus constructions using the common degree-one circle
coordinate.  Their product areas converge to their base areas as
$\lambda\downarrow0$, uniformly over the finite cover, so one
$\lambda_{\rm net}>0$ makes them all smaller than $\mu$.

The continuous curve-shrinking family is supplied by
\ref{thm:controlled-curve-shrinking-existence}.  Projection to the circle
factor has degree $1$ on every ramp and is preserved by the curve homotopy.
The projection-length inequality in
\ref{lem:ramp-annulus-nondegeneration} therefore gives
$L(F\circ\alpha)\ge L(S^1_\lambda)=\lambda$ for every essential domain
circle $\alpha$, simultaneously for all the comparisons just chosen.  The
last assertion follows by taking the minimum of finitely many positive
numbers.  This is the quantifier order in the source: choose the annular net,
then impose the finitely many pointwise small-$\lambda$ requirements, and
only then fix one common positive $\lambda$.
\end{proof}

\begin{lemma}[Imported annular approximation and variation interface]
\label{lem:annular-approximation-comparison}
\dcref{MT:ch19:19.15}
\source{morgan-tian}
\uses{thm:annular-douglas-minimizer}
\notready
For a smooth pair of degree-one ramp flows in dimension at least three, the
pair can be approximated in $C^2$ by disjoint real-analytic pairs in the same
free homotopy class.  The two boundary pairs cobound thin annuli whose areas
tend to zero, and the corresponding annular infima and forward difference
quotients converge in both directions.  This is the source's third-curve and
analytic-regularization interface.
\end{lemma}

\begin{proof}
This is the analytic-regularization and third-curve interface in
Morgan--Tian, Chapter~19, Proposition~19.15.  Approximate the smooth ramp
pair in $C^2$ by a disjoint real-analytic pair in the same free homotopy
class.  Short geodesic collars between the original and approximating
boundaries have area tending to zero, so gluing them to admissible annuli
gives the two inequalities for the corresponding infima.  Apply the same
collar comparison to the forward difference quotients and then take the
analytic approximants to the smooth pair.  The degree-one circle projection
keeps the nondegeneration lower bound uniform throughout this passage.  The
attainment and branch hypotheses used for each analytic minimizer are exactly
those stated in \ref{thm:annular-douglas-minimizer}.
\end{proof}

\begin{lemma}[Variation of the annular area infimum]
\label{lem:annular-area-infimum-variation}
\dcref{MT:ch19:19.15}
\source{morgan-tian}
\uses{def:ramp-curve,thm:annular-douglas-minimizer,
  lem:ramp-annulus-nondegeneration,lem:smooth-annular-minimizer-bridge,
  lem:annular-approximation-comparison}
\notready
Let $n=\dim M\ge3$, and let $c_0(t),c_1(t)$ be degree-one ramp
curve-shrinking flows in $(M,g(t))\times S^1_\lambda$ in the same free
homotopy class.  Assume at every time under consideration that the positive
essential-loop lower bound of
\ref{lem:ramp-annulus-nondegeneration} holds.  The theorem
\ref{thm:annular-douglas-minimizer} then supplies the minimizing annulus.  If
$\mu(t)$ is the infimum of the areas of spanning annuli,
then $\mu$ is continuous and
\[
 D^+\mu(t)\le(2n-1)\max_M|\operatorname{Rm}_{g(t)}|\,\mu(t).
\]
Consequently $\mu(t)\le e^{C(t-s)}\mu(s)$ on a compact time interval, with
$C$ independent of $\lambda$.
\end{lemma}

\begin{proof}
 The smooth-flow bridge \ref{lem:smooth-annular-minimizer-bridge}, together with
\ref{lem:ramp-annulus-nondegeneration}, supplies a minimizing annulus and the
uniform positive length bound for every essential domain circle.  First take
disjoint boundary curves.  Move the minimizer's boundary by the two curve
flows; interior normal variation vanishes by minimality.  Metric and boundary
variation give
\[
 \left.\frac d{ds}\right|_{s=t}\operatorname{Area}(A_s)
 =-\int_A\operatorname{tr}_{TA}\operatorname{Ric}\,da
   -\int_{\partial A}k_g\,ds.
\]
Apply Gauss--Bonnet on parallel subannuli avoiding the finite branch set and
let them tend to the boundary.  If the interior branch orders are $m_j$, the
result is
\[
 -\int_{\partial A}k_g\,ds
 =\int_AK_A\,da-2\pi\sum_j(m_j-1)
 \le\int_AK_A\,da.
\]
Minimality gives $K_A\le\sec_{M\times S^1_\lambda}(TA)$, and
$|\operatorname{tr}_{TA}\operatorname{Ric}|\le2(n-1)\max|\operatorname{Rm}|$.  The circle is flat, so
the constants are independent of $\lambda$.  This proves the upper
derivative estimate.

In dimension $n\ge3$, apply the imported approximation interface
\ref{lem:annular-approximation-comparison} to pass from the disjoint analytic
calculation to the stated smooth pair.  Its two-sided thin-annulus comparison
gives continuity and convergence of the forward difference quotients.  Scalar
integrating-factor comparison gives the exponential estimate.  All attainment
and branch-regularity inputs are supplied by
\ref{thm:annular-douglas-minimizer} from the displayed nondegeneration bound.
\end{proof}

\begin{theorem}[Imported small-annulus length transfer]
\label{thm:small-annulus-length-transfer}
\dcref{MT:ch19:19.31,MT:ch19:19.35}
\source{morgan-tian}
\uses{def:ramp-curve,thm:annular-douglas-minimizer,
  lem:ramp-annulus-nondegeneration,
  lem:smooth-annular-minimizer-bridge,
  lem:annular-approximation-comparison}
\notready
There is $\delta_0>0$ with the following property.  For every $r>0$ and every
ambient sectional-curvature upper bound $K$, there is $\mu_0>0$ such that if
two ramps $\gamma,\widehat\gamma$ in
$M\times S^1_\lambda$ cobound an admissible annulus of area below $\mu_0$, and
\[
 L(\gamma)\ge r,
 \qquad
 \int_I|\mathbf k_\gamma|\,ds<\delta_0
\]
for every length-$r$ subarc $I\subset\gamma$, where
$\mathbf k_\gamma$ is the ambient curvature vector of the unit-speed ramp
$\gamma$, then
\[
 L(\widehat\gamma)\ge\frac34L(\gamma).
\]
The constants are uniform in $\lambda$.  For each annular comparison to
which this contract is applied, the spanning annulus is required to satisfy
the nondegeneration bound in
\ref{lem:ramp-annulus-nondegeneration}.  The analytic/disjoint-boundary
approximation, Douglas minimizer, branch-point metric regularization, and
limiting passage are those of
\ref{lem:smooth-annular-minimizer-bridge} and
\ref{lem:annular-approximation-comparison}; after that passage the preceding
Douglas theorem supplies the minimizer whenever the proof uses one.
\end{theorem}

\begin{proof}
This is the intrinsic annulus estimate of Morgan--Tian, Chapter~19,
Propositions~19.31 and~19.35.  Assume the asserted conclusion fails for a
sequence of admissible annuli with area tending to zero.  First use
\ref{lem:annular-approximation-comparison} and
\ref{lem:smooth-annular-minimizer-bridge} to perturb to disjoint analytic
boundaries and replace each competitor by its Douglas minimizer.  Remove the
finitely many interior branch points by the branch-point metric
regularization in that bridge.  Along the boundary,
$k_{\rm geod}=\langle\mathbf k_\gamma,n\rangle$, hence
$|k_{\rm geod}|\le|\mathbf k_\gamma|$.  Use the lower bound on every
essential domain circle and apply Gauss--Bonnet together with the boundary
geodesic-curvature and focal-radius estimates.  The small area forces the two
boundary curves to remain within the controlled collar for every length-$r$
subarc; integrating the turning bound then gives
$L(\widehat\gamma)\ge(1-\delta_0')L(\gamma)$, with
$\delta_0'$ chosen below $1/4$.  This contradiction yields the displayed
$3/4$ constant.  For general smooth ramps, the same approximation/bridge
perturbs the boundaries to disjoint real-analytic ramps, applies the estimate
to their Douglas minimizers, and passes to the limit through the thin collar
annuli.  The estimates use only
the intrinsic annulus metric and the ambient sectional-curvature upper bound,
so the constants are uniform in the flat circle scale $\lambda$.  The
nondegeneration hypothesis is explicit, and Douglas attainment is supplied
by \ref{thm:annular-douglas-minimizer} after the stated approximation.
\end{proof}

\begin{lemma}[Single-ramp terminal dichotomy]
\label{lem:single-ramp-terminal-dichotomy}
\dcref{MT:ch18:18.24}
\source{morgan-tian}
\uses{def:loop-width-comparison-ode,
  thm:least-disk-area-variation-interface,
  thm:controlled-curve-shrinking-existence,
  lem:curve-flow-length-total-curvature-evolution,
  thm:curve-flow-length-total-curvature-bounds}
\notready
In the setting of
\ref{thm:controlled-curve-shrinking-existence}, let
$p_1\colon M\times S^1_\lambda\to M$ be product projection and fix
$\zeta>0$.  There is
$B>(t_1-t_0)^{-1}$, depending only on $\zeta$, the approximating family, and
the ambient Ricci-flow bounds, such that, with $t_2=t_1-B^{-1}$, the following
holds.  For every $c\in S^2$ there is $\lambda_c>0$ such that for every
$0<\lambda<\lambda_c$ at least one of these alternatives holds:
\begin{enumerate}
\item
\[
 \mathcal A_{g(t_1)}
   \bigl(p_1\widetilde\Gamma_c^\lambda(t_1)\bigr)
 <w_{\mathcal A_{g(t_0)}(\widetilde\Gamma(c)),t_0}(t_1)+\zeta/2;
\]
\item
\[
 L\bigl(\widetilde\Gamma_c^\lambda(t)\bigr)<\zeta/2
 \qquad\text{for every }t\in[t_2,t_1].
\]
\end{enumerate}
The number $B$ is common to the whole compact family; only the preliminary
threshold $\lambda_c$ is pointwise.  More precisely, the length-growth and
stopping estimates furnish constants $C_2,C_5,\delta'>0$ (depending only on
the same family and ambient bounds), and $B$ is chosen so that
\[
 B\ge \frac{3C_5}{\delta'},\qquad
 B\ge \frac{3e^{C_2(t_1-t_0)}}{\zeta},\qquad
 B>\frac{C_2}{\log 4-\log 3}.
\]
The last inequality is the quantitative implication
$e^{C_2/B}<4/3$ used in the short-transfer argument.
\end{lemma}

\begin{proof}
This is the stronger one-curve assertion called
\texttt{xistronger} in the proof of Morgan--Tian Proposition~18.24
(\texttt{energy1.tex}, lines~2510--2522).  The curve-flow length and
total-curvature estimates bound the number and total size of the bad time
intervals independently of $\lambda$.  On the complementary intervals,
the least-disk variation inequality and scalar comparison give the first
alternative.  If that comparison cannot be continued through the final
interval, the stopping estimates give the second alternative on the entire
interval $[t_2,t_1]$.  Choosing $B$ to satisfy the three displayed
inequalities makes it independent of $c$ and supplies the stopping and
backward-growth margins; the limiting argument in $\lambda\downarrow0$ supplies
the pointwise $\lambda_c$.
\end{proof}

\begin{lemma}[A short net ramp forces a short neighboring terminal loop]
\label{lem:finite-net-short-terminal-transfer}
\dcref{MT:ch18:18.24,MT:ch19:19.31,MT:ch19:19.35}
\source{morgan-tian}
\uses{lem:single-ramp-terminal-dichotomy,
  lem:annular-area-infimum-variation,
  thm:small-annulus-length-transfer,
  lem:curve-flow-length-total-curvature-evolution,
  thm:curve-flow-length-total-curvature-bounds}
\notready
Fix the compact ramp family, $\zeta$, $B$, and
$t_2=t_1-B^{-1}$ from
\ref{lem:single-ramp-terminal-dichotomy}, with $p_1$ denoting the same product
projection.  There is $\mu_{\rm short}>0$,
independent of $0<\lambda<1$, with the following property.  Suppose
$c,\widehat c\in S^2$ and the initial ramps
$\widetilde\Gamma_c^\lambda(t_0)$ and
$\widetilde\Gamma_{\widehat c}^\lambda(t_0)$ cobound an annulus of area less
than $\mu_{\rm short}$.  If
\[
 L\bigl(\widetilde\Gamma_{\widehat c}^\lambda(t)\bigr)<\zeta/2
 \qquad(t_2\le t\le t_1),
\]
then
\[
 L_{g(t_1)}\bigl(p_1\widetilde\Gamma_c^\lambda(t_1)\bigr)<\zeta.
\]
\end{lemma}

\begin{proof}
This is the claim \texttt{mulength} in the proof of Morgan--Tian
Proposition~18.24 (\texttt{energy1.tex}, lines~2893--2956).  If the projected
terminal loop at $c$ had length at least $\zeta$, exponential length growth,
with the defining choice of $B$, would keep its ramp length at least
$3\zeta/4$ throughout $[t_2,t_1]$.  Integrating the length evolution gives a
time $t'$ in that interval with a uniform $L^2$ curvature bound.  After
choosing a sufficiently short subarc scale $r$, Cauchy--Schwarz supplies the
small-turning hypothesis of
\ref{thm:small-annulus-length-transfer}.  By
\ref{lem:annular-area-infimum-variation}, an initial annulus smaller than a
sufficiently small $\mu_{\rm short}$ remains below the length-transfer
threshold at $t'$.  That theorem then gives
\[
 L\bigl(\widetilde\Gamma_{\widehat c}^\lambda(t')\bigr)
 \ge\frac34L\bigl(\widetilde\Gamma_c^\lambda(t')\bigr)
 \ge\frac9{16}\zeta>\frac12\zeta,
\]
contrary to the assumed persistent shortness of the net ramp.  All ambient
and curve-flow constants are uniform in $\lambda$, so the same
$\mu_{\rm short}$ works for the family.
\end{proof}


\begin{theorem}[Imported controlled deformation of a loop family]
\label{thm:controlled-loop-family-deformation}
\dcref{MT:ch18:18.24}
\source{morgan-tian}
\uses{def:loop-family-width,def:loop-width-comparison-ode,
  lem:loop-width-ode-formula,lem:least-spanning-area-continuity,
  thm:least-disk-area-variation-interface,
  thm:annular-douglas-minimizer,
  lem:ramp-annulus-nondegeneration,
  def:finite-small-annulus-net,
  lem:finite-net-ramp-nondegeneration,
  lem:annular-area-infimum-variation,
  lem:annular-approximation-comparison,
  thm:small-annulus-length-transfer,
  lem:single-ramp-terminal-dichotomy,
  lem:finite-net-short-terminal-transfer,
  thm:controlled-curve-shrinking-existence}
\notready
Let $g(t)$ be a Ricci flow on a compact three-manifold for
$t\in[t_0,t_1]$, let $\Gamma\colon S^2\to\Lambda_0M$ be continuous, and let
$\zeta>0$.  There is a continuous family
$\widetilde\Gamma(t)\colon S^2\to\Lambda_0M$ such that
the maps $\widetilde\Gamma(t_0)$ and $\Gamma$ are homotopic as maps
$S^2\to\Lambda_0M$, and the path $u\mapsto\widetilde\Gamma(u)$ is a homotopy
of maps from $\widetilde\Gamma(t_0)$ to $\widetilde\Gamma(t)$ for every
$t\in[t_0,t_1]$.  In particular, whenever $\pi_2(M)=0$, all these maps have
the same free $\pi_1(M)$-orbit under the established loop-space
$\pi_2=0$ identification.  For every $c\in S^2$,
\[
 |\mathcal A_{g(t_0)}(\widetilde\Gamma(t_0)(c))
   -\mathcal A_{g(t_0)}(\Gamma(c))|<\zeta,
\]
and at least one of the following holds:
\begin{enumerate}
\item $L_{g(t_1)}(\widetilde\Gamma(t_1)(c))<\zeta$;
\item
\[
 \mathcal A_{g(t_1)}(\widetilde\Gamma(t_1)(c))
 \le w_{\mathcal A_{g(t_0)}(\widetilde\Gamma(t_0)(c)),t_0}(t_1)+\zeta.
\]
\end{enumerate}
More precisely, the construction first fixes the common terminal interval
from \ref{lem:single-ramp-terminal-dichotomy}, then a positive annular-area
threshold $\mu$, then a finite $\mu/2$-annular net, and only afterward one
common scale $\lambda>0$ below the finitely many pointwise thresholds and
$\lambda_{\rm net}$ of
\ref{lem:finite-net-ramp-nondegeneration}.  At that fixed scale every net
comparison starts with annular area below $\mu$ and has the positive
nondegeneration constant $\rho_\lambda=\lambda$ throughout the time
interval.  Douglas attainment and the corresponding branch/collar
regularity are then supplied by
\ref{thm:annular-douglas-minimizer}; no uniform positive lower bound as
$\lambda\downarrow0$ is asserted.
\end{theorem}

\begin{proof}
This is the finite-net deformation construction in Morgan--Tian,
Chapter~18, Proposition~18.24.  First use
\ref{thm:controlled-curve-shrinking-existence} with approximation error much
smaller than $\zeta$, and call the resulting initial smooth family
$\Gamma^{\rm app}$.  Apply
\ref{lem:single-ramp-terminal-dichotomy} with a correspondingly smaller error
to obtain one $B$ and $t_2=t_1-B^{-1}$ for the whole family.  Let
$\mu_{\rm short}$ be the threshold in
\ref{lem:finite-net-short-terminal-transfer}.  By
\ref{lem:annular-area-infimum-variation} and the initial-area sensitivity in
\ref{lem:loop-width-ode-formula}, choose $0<\mu<\mu_{\rm short}$ so small that
an initial annular infimum below $\mu$ contributes less than the remaining
error both at time $t_1$ and in the terminal scalar comparison.

Now apply \ref{lem:finite-net-ramp-nondegeneration} with this $\mu$.  It gives
a finite net ${\mathcal S}$ and $\lambda_{\rm net}$.  For every
$\widehat c\in{\mathcal S}$, the single-ramp dichotomy gives a positive
$\lambda_{\widehat c}$.  The common-scale clause of the finite-net lemma
allows one choice
\[
 0<\lambda<\min\{\lambda_{\rm net},
        \lambda_{\widehat c}:\widehat c\in{\mathcal S}\}.
\]
Fix $c\in S^2$ and choose $\widehat c\in{\mathcal S}$ with initial ramp
annulus of area below $\mu$.  The annular variation theorem keeps its
infimum below the prescribed terminal error.  If the area alternative of
\ref{lem:single-ramp-terminal-dichotomy} holds at $\widehat c$, glue a
terminal comparison annulus to a spanning disk and use
\ref{lem:least-spanning-area-continuity} together with the scalar ODE
sensitivity to obtain the displayed area alternative at $c$.  If instead
the net ramp is shorter than $\zeta/2$ throughout $[t_2,t_1]$, apply
\ref{lem:finite-net-short-terminal-transfer}; this gives the displayed
terminal length alternative at $c$.  This is exactly the
$\mu$-net/common-$\lambda$ order of the source proof
(\texttt{energy1.tex}, lines~2960--2981).

Finally set $\widetilde\Gamma(t)$ equal to the projection of this one
common-$\lambda$ ramp family.  The parametric curve flow gives a continuous
map $S^2\times[t_0,t_1]\to\Lambda_0M$, and its time-$t_0$ member is freely
homotopic as a map to $\Gamma$ by the approximation contract.  Hence the
time path supplies the asserted homotopy of maps; when $\pi_2(M)=0$, the
loop-space identification turns it into preservation of the same free-loop
orbit.  The initial polygonal and
thin-collar errors were chosen below $\zeta$, which proves the initial area
estimate and completes both terminal alternatives.
\end{proof}


\section{\texorpdfstring{$\pi_3$}{pi3}-Width Evolution Through Surgery}

\begin{lemma}[Coherent marked-point transport along a component path]
\label{lem:coherent-marked-component-path}
\dcref{MT:ch18:18.18,MT:ch15:15.12}
\source{morgan-tian}
\uses{def:path-of-components}
\notready
Let $\mathcal X$ be a component path with finitely many singular times.
There are marked points $x(t)\in\mathcal X(t)$, continuous on each regular
interval, and paths in the continuing surgery components joining
$F_t(x(t))$ to $x(T)$ at every singular time $T$.  These choices identify
the based homotopy groups across regular intervals and make each surgery map
an induced map after the displayed basepoint change.  A different choice
changes the based identification only by the standard $\pi_1$-action; hence
the induced map on the orbit set of $\pi_3$ is canonical.
\end{lemma}

\begin{proof}
Choose the marked point in the compact smooth part away from all finitely many
neck and cap regions.  The time vector field transports it on each regular
interval.  At a surgery, the continuing component is path connected and the
Lipschitz pinch map is defined on that compact core, so join its image to the
next marked point by a path in the core.  Basepoint-change isomorphisms along
two such paths differ by conjugation by a loop, which is exactly the natural
$\pi_1$-action on $\pi_3$.  This proves both coherence and independence on
orbit classes.
\end{proof}

\begin{theorem}[Transport of a nonzero third-homotopy class]
\label{thm:surgery-pi3-class-transport}
\dcref{MT:ch18:18.18,MT:ch15:15.12}
\source{morgan-tian}
  \uses{def:path-of-components,thm:trivial-surgery-component-map,
  thm:finite-interval-surgery-bound,
  lem:nullhomotopic-embedded-sphere-separates,
  lem:coherent-marked-component-path,
  thm:loop-space-pi3-identification}
\notready
Let $\mathcal X$ be a component path on $[t_0,t_1]$ with only finitely many
surgery times, with
$\pi_2(\mathcal X(t_0))=0$ and
$\pi_3(\mathcal X(t_0))\cong\mathbb Z$, and let
$0\ne\xi(t_0)\in\pi_3(\mathcal X(t_0))$.  Choose the coherent marked
points of \ref{lem:coherent-marked-component-path}.  Then $\pi_2(\mathcal
X(t))=0$
for every $t$, and there is a compatible family of nonzero classes
$\xi(t)\in\pi_3(\mathcal X(t))$.  At a surgery time $T$, the maps
$F_t\colon\mathcal X(t)\to\mathcal X(T)$ of
\ref{thm:trivial-surgery-component-map}, for $t<T$ sufficiently close to $T$,
satisfy
\[
 (F_t)_*[\xi(t)]=[\xi(T)]
\]
in the $\pi_1$-orbit sets; after the coherent marked-point choices this
equality may be read as an equality of based representatives.
\end{theorem}

\begin{proof}
At a regular time, the marked-point transport of
\ref{lem:coherent-marked-component-path} identifies the homotopy groups and
the orbit of the class.  At a surgery time, every cutting sphere is
null-homotopic because
$\pi_2=0$, and it is separating by
\ref{lem:nullhomotopic-embedded-sphere-separates}.  Apply
\ref{thm:trivial-surgery-component-map}.  It preserves the
vanishing of $\pi_2$, and its map on the relevant infinite cyclic third
homotopy groups is injective.  Hence the induced orbit class
$(F_t)_*[\xi(t)]$ is nonzero.  Induction over the surgery times gives the
compatible orbit family; the based representatives are recovered from the
coherent marked points.
\end{proof}

For the transported class in \ref{thm:surgery-pi3-class-transport}, write
$W_\xi(t)=W_{g(t)}(\xi(t))$.

\begin{lemma}[Continuity of loop-family width at regular times]
\label{lem:loop-width-regular-continuity}
\dcref{MT:ch18:18.21}
\source{morgan-tian}
\uses{def:loop-family-width,thm:surgery-pi3-class-transport}
\notready
The function $W_\xi$ is continuous at every nonsurgery time.
\end{lemma}

\begin{proof}
Use the time vector field to identify a closed surgery-free neighborhood of
$t$ with a fixed compact manifold.  The pulled-back metrics satisfy
\[
 c(s)^{-1}g(t)\le g(s)\le c(s)g(t),
 \qquad c(s)\longrightarrow1.
\]
Every Lipschitz disk then has its area multiplied by a factor between
$c(s)^{-1}$ and $c(s)$.  The same bounds pass first to the infimum defining
$\mathcal A$, then to the maximum over a family, and finally to the infimum
over representatives of the transported class.  Thus
\[
 c(s)^{-1}W_\xi(t)\le W_\xi(s)\le c(s)W_\xi(t),
\]
which proves continuity.
\end{proof}

\begin{lemma}[Loop-family width at a surgery time]
\label{lem:loop-width-surgery-lower-semicontinuity}
\dcref{MT:ch18:18.22,MT:ch15:15.12}
\source{morgan-tian}
\uses{def:loop-family-width,thm:surgery-pi3-class-transport,
  thm:trivial-surgery-component-map}
\notready
At a surgery time $T$,
\[
 W_\xi(T)\le\liminf_{t\uparrow T}W_\xi(t).
\]
\end{lemma}

\begin{proof}
For every $\eta>0$ and all $t<T$ sufficiently close to $T$,
\ref{thm:trivial-surgery-component-map} gives a map $F_t$ with Lipschitz
constant at most $1+\eta$.  By
\ref{thm:surgery-pi3-class-transport}, it sends $\xi(t)$ to $\xi(T)$.  If a
Lipschitz disk spans $\gamma$, then $F_t$ composed with that disk spans
$F_t\circ\gamma$ and has area at most $(1+\eta)^2$ times the original area.
Apply this pointwise to a representative family and take both infima to get
\[
 W_\xi(T)\le(1+\eta)^2W_\xi(t).
\]
Take a sequence realizing the left liminf and then let $\eta\downarrow0$.
\end{proof}

\begin{lemma}[Short loops bound small-area disks]
\label{lem:short-loops-small-disk}
\dcref{MT:ch18:18.28}
\source{morgan-tian}
\notready
Let $(X,g)$ be compact.  For every $\eta>0$ there is
$0<\zeta<\eta/2$ such that every $C^1$ loop of length less than $\zeta$
bounds a Lipschitz disk of area less than $\eta$.
\end{lemma}

\begin{proof}
Choose $r>0$ below the normal radius uniformly over $X$.  If a loop $\gamma$
has length $\ell<r/2$, put $p=\gamma(1)$ and write
$v(\theta)=\exp_p^{-1}(\gamma(\theta))$.  The radial filling
\[
 F(s,\theta)=\exp_p(sv(\theta))
\]
is defined on the disk.  On the compact radius-$r$ normal bundle the
differential of the exponential map is uniformly bounded.  Hence
$|\partial_sF|\le C\ell$ and
$|\partial_\theta F|\le C|v'(\theta)|$, so
\[
 \operatorname{Area}(F)\le C^2\ell\int_{S^1}|v'|\,d\theta\le C'\ell^2.
\]
Choose
$\zeta<\min\{\eta/2,r/2,\sqrt{\eta/C'}\}$.  Then every loop of length below
$\zeta$ has the required filling.
\end{proof}

\begin{lemma}[Short loop families are trivial]
\label{lem:short-loop-families-trivial}
\dcref{MT:ch18:18.27}
\source{morgan-tian}
\uses{thm:loop-space-pi3-identification}
\notready
Let $(X,g)$ be compact with $\pi_2(X)=0$.  There is $\zeta>0$ such that any
continuous family $\Gamma\colon S^2\to\Lambda_0X$ all of whose loops have
length less than $\zeta$ represents the zero free-homotopy orbit, and hence
the zero class after any based choice, under
$\pi_2(\Lambda_0X,*)\cong\pi_3(X)$.
\end{lemma}

\begin{proof}
Choose $\zeta$ below the normal radius.  For each $c\in S^2$, join
$\Gamma(c)(\theta)$ to $\Gamma(c)(1)$ by the unique short geodesic.  Smooth
dependence of the geodesic on its endpoints gives a continuous family of
radial disks.  Shrinking through those disks homotopes $\Gamma$ to the family
of constant loops $c\mapsto\Gamma(c)(1)$.  This remaining family is a map
$S^2\to X$, and is null-homotopic because $\pi_2(X)=0$.  The identification
in \ref{thm:loop-space-pi3-identification} gives the conclusion.
\end{proof}

\begin{theorem}[Forward inequality for the $\pi_3$ width]
\label{thm:loop-width-forward-inequality}
\dcref{MT:ch18:18.18}
\source{morgan-tian}
\uses{def:path-of-components,def:loop-family-width,
  def:loop-width-comparison-ode,
  lem:loop-width-ode-formula,thm:controlled-loop-family-deformation,
  def:finite-small-annulus-net,
  lem:finite-net-ramp-nondegeneration,
  lem:single-ramp-terminal-dichotomy,
  lem:finite-net-short-terminal-transfer,
  lem:short-loops-small-disk,lem:short-loop-families-trivial,
  lem:loop-width-regular-continuity,
  lem:loop-width-surgery-lower-semicontinuity,
  thm:surgery-pi3-class-transport}
\notready
Let $\mathcal X$ be a component path of a controlled surgery flow on a compact
interval $[t_0,t_1]$ containing only finitely many surgery times.  Suppose
$\pi_2(\mathcal X(t_0))=0$ and
$\pi_3(\mathcal X(t_0))\cong\mathbb Z$, and let
$0\ne\xi(t_0)\in\pi_3(\mathcal X(t_0))$.  Let
$\xi(t)$ be the compatible family of nonzero classes supplied by
\ref{thm:surgery-pi3-class-transport}; in particular
$\pi_2(\mathcal X(t))=0$ for every $t$.  Then
$W_\xi$ is continuous at nonsurgery times, obeys
\[
 W_\xi(T)\le\liminf_{t\uparrow T}W_\xi(t)
\]
at surgery times, and satisfies
\[
 D^+W_\xi(t)
 \le-2\pi-\frac12R_{\min}(t)W_\xi(t)
\]
on the immediate right of every time in $[t_0,t_1)$ for which the component
path survives.
\end{theorem}

\begin{proof}
The transport and continuity assertions are
\ref{thm:surgery-pi3-class-transport},
\ref{lem:loop-width-regular-continuity}, and
\ref{lem:loop-width-surgery-lower-semicontinuity}.  Fix a surgery-free
interval $[s,t]$ and let $w$ solve
\ref{def:loop-width-comparison-ode} with $w(s)=W_\xi(s)$.  We prove
$W_\xi(t)\le w(t)$ whenever the component survives to $t$.

First suppose $w(t)\ge0$.  Given $\eta>0$, choose a representative family
with width below $W_\xi(s)+\varepsilon$.  Use
\ref{lem:loop-width-ode-formula} and then choose the deformation error
$\zeta$ so small that changing any initial disk area by at most
$2\zeta$ changes its terminal ODE value by less than $\eta/2$.  Also choose
$\zeta$ by \ref{lem:short-loops-small-disk} so that every terminal loop of
length below $\zeta$ spans a disk of area below $\eta$.  Apply
\ref{thm:controlled-loop-family-deformation}.  In its area alternative the
terminal area is below $w(t)+\eta$; in its short alternative it is below
$\eta\le w(t)+\eta$.  The explicit net clause of the deformation theorem,
supplied by \ref{lem:finite-net-ramp-nondegeneration}, makes one common
$\lambda$ satisfy the annular-area and single-ramp thresholds; its short
branch is transferred to every parameter by
\ref{lem:finite-net-short-terminal-transfer}.  Its homotopy clause, together
with $\pi_2(\mathcal X(t))=0$, says that the terminal family still represents
the free orbit $[\xi(t)]$.  Hence it is an
admissible family in the definition of $W_\xi(t)$ and has width at most
$w(t)+\eta$.  Let $\eta,\varepsilon\downarrow0$.

If $w(t)<0$, choose $\eta>0$ with $w(t)+\eta<0$.  Choose
$\zeta>0$ below $\eta/2$ and the threshold in
\ref{lem:short-loop-families-trivial}.  Apply
$\ref{lem:loop-width-ode-formula}$ and then choose
$\varepsilon>0$ (shrinking $\zeta$ further if necessary) so that the ODE
sensitivity estimate sends every initial area below
$W_\xi(s)+\varepsilon+\zeta$ to a terminal value below
$w(t)+\eta/2$.  Apply
\ref{thm:controlled-loop-family-deformation} to a representative of width
below $W_\xi(s)+\varepsilon$.  Its area alternative would then give a
terminal area below $w(t)+\eta/2+\zeta<w(t)+\eta<0$, hence a negative
spanning area, and is impossible.  Thus every terminal loop is
short, so the terminal family is homotopically trivial.  By the explicit
orbit-preservation output of
\ref{thm:controlled-loop-family-deformation} and
$\pi_2(\mathcal X(t))=0$, that terminal orbit is $[\xi(t)]$.  This contradicts
transport of the nonzero class.  Consequently
the component cannot survive to a time at which the comparison solution is
negative.

Apply the comparison just proved on arbitrarily short regular intervals.
Together with regular-time continuity it yields the displayed upper forward
derivative.  Starting immediately after a surgery gives the same right-hand
inequality there; the separate left lower-semicontinuity statement accounts
for the jump from the preceding interval.
\end{proof}

\section{Finite Extinction}

\begin{lemma}[Imported prime-factor contract for a two-sided projective plane]
\label{lem:projective-plane-prime-factor-contract}
\dcref{MT:ch18:18.1}
\source{morgan-tian}
\notready
Let $M$ be a closed connected three-manifold whose fundamental group is a
free product of finite groups and infinite cyclic groups, and let $P\subset M$
be an embedded two-sided $\mathbb{RP}^2$ with trivial normal bundle.  Under
the prime/connected-sum decomposition used for this free product, $P$ can be
isotoped disjoint from the decomposition spheres and hence into one prime
factor.  Its inclusion is $\pi _1$-injective, and because it is two-sided and
nonseparating in that factor, the factor's first homology surjects onto
$\mathbb Z$.  This is the Hempel/prime-decomposition topology input in the
Morgan--Tian proof.
\end{lemma}

\begin{proof}
This is the prime-decomposition input used in Morgan--Tian, Chapter~18,
Theorem~18.1.  Put $P$ in general position with respect to the finite family
of decomposition spheres.  The innermost-circle argument removes all
intersections, so the plane lies in a single prime factor.  Its two-sided
complement and the induced orientation double cover give a nonseparating loop
in that factor; the inclusion of $\pi_1(\mathbb{RP}^2)=\mathbb Z/2$ is
injective, while the dual loop gives a surjection of the factor's first
homology onto $\mathbb Z$.  The Mayer--Vietoris step is the homology
calculation recorded in Morgan--Tian's Chapter~18 proof; the prime-factor and
Hempel classification are imported exactly as Morgan--Tian state them, with
the normal-bundle and free-product hypotheses explicit above.  The registered
Morgan--Tian source is the authoritative import used here, so no unregistered
page-level Hatcher anchor is being asserted.
\end{proof}

\begin{theorem}[The free-product hypothesis excludes a two-sided projective
plane]
\label{thm:free-product-excludes-two-sided-projective-plane}
\dcref{MT:ch18:18.1}
\source{morgan-tian}
\uses{lem:projective-plane-prime-factor-contract}
\notready
Let $M$ be a closed connected three-manifold whose
fundamental group is a free product of finite groups and infinite cyclic
groups.  Then $M$ contains no embedded $\mathbb{RP}^2$ with trivial normal
bundle.
\end{theorem}

\begin{proof}
Suppose such a projective plane $P$ exists.  Apply
\ref{lem:projective-plane-prime-factor-contract}.  The prime factor
containing $P$ has an infinite fundamental group because its first homology
surjects onto $\mathbb Z$, but it also contains the injective order-two
subgroup coming from $\pi _1(P)$.  It is therefore neither finite nor
infinite cyclic, contradicting the assumed free-product factor list.  This
is the topological argument in Morgan--Tian 18.1; the prime-factor and
Hempel assertions are the explicit source-backed import in
\ref{lem:projective-plane-prime-factor-contract}.
\end{proof}

\begin{theorem}[Finite-time extinction]
\label{thm:finite-time-extinction}
\dcref{MT:ch18:18.1}
\source{morgan-tian}
\uses{def:path-of-components,
  def:controlled-ricci-flow-with-surgery,
  def:ricci-flow-with-cutoff,
  thm:all-time-controlled-surgery-flow,
  thm:finite-interval-surgery-bound,
  thm:eventual-pi2-vanishing,
  thm:late-component-nontrivial-pi3,
  thm:loop-width-forward-inequality,
  lem:extinction-scalar-curvature-lower-bound,
  lem:forward-dini-comparison-with-downward-jumps}
\notready
Let $(M,g_0)$ be a compact connected normalized Riemannian three-manifold
whose fundamental group is a free product of finite groups and infinite
cyclic groups and which contains no embedded $\mathbb{RP}^2$ with trivial
normal bundle.  Fix an all-time three-dimensional generalized Ricci flow
$({\mathcal M},G)$ starting from $(M,g_0)$ which has the
Kleiner--Lott $(r,\delta)$-cutoff, satisfies Assumptions (1)--(7), has
curvature pinched toward positive, satisfies the stagewise strong canonical
neighborhood and noncollapsing estimates, and has only finitely many surgery
times on each bounded interval, and which is one of the all-time flows
supplied by \ref{thm:all-time-controlled-surgery-flow}.  Then this fixed flow has a
finite extinction time: there is $T_*<\infty$ such that
$M_t=\varnothing$ for every $t\ge T_*$.
\end{theorem}

\begin{proof}
Choose $T_1$ from \ref{thm:eventual-pi2-vanishing}.  If a component path
$\mathcal X$ survives on $[T_1,T]$, then
\ref{thm:late-component-nontrivial-pi3} supplies a nonzero class
$\xi\in\pi_3(\mathcal X(T_1))$.

By \ref{thm:loop-width-forward-inequality} and
\ref{lem:extinction-scalar-curvature-lower-bound},
\[
 D^+W_\xi(t)\le-2\pi+\frac3{1+4t}W_\xi(t).
\]
Let $w(T_1)=W_\xi(T_1)$ and
$w'=-2\pi+3w/(1+4t)$.  Direct integration gives the consistent formula
\[
 w(t)=W_\xi(T_1)
       \left(\frac{1+4t}{1+4T_1}\right)^{3/4}
 +2\pi(1+4T_1)^{1/4}(1+4t)^{3/4}-2\pi(1+4t).
\]
The finite-partition comparison lemma
\ref{lem:forward-dini-comparison-with-downward-jumps}, applied to the finite
list of surgery times on each bounded interval, combines the regular-interval
inequality with the downward surgery liminf and gives
$0\le W_\xi(t)\le w(t)$ while the path survives.  The negative linear
term in $w$ dominates its $t^{3/4}$ terms, so $w$ becomes negative after a
finite time depending only on $T_1$ and $W_\xi(T_1)$.  Thus this component
path must end before that time.

The compact time slice $M_{T_1}$ has finitely many components.  Every
component at a later time has an ancestor among them, represented by a path of
components in the sense of \ref{def:path-of-components}; bounded-interval
surgery finiteness ensures that the transport induction along such a path is
finite.  Choose one generator for each nonzero cyclic third-homotopy group at
time $T_1$.  The preceding comparison bound depends only on that component and
generator, not on which descendant is selected at a surgery.  Taking the
maximum of the finitely many lifetime bounds gives $T_*<\infty$ after which
no component remains.  Once a time slice is empty, every later slice is empty
by the definition of the surgery flow.
\end{proof}
